\documentclass[11pt]{article}
\usepackage[margin=1in]{geometry}
\usepackage{amsmath,amssymb,amsthm,hyperref}
\newtheorem{lemma}{Lemma}
\newtheorem{proposition}[lemma]{Proposition}
\newcommand{\R}{\mathbb R}
\newcommand{\diag}{\operatorname{diag}}
\title{A fixed completion on the compact Bob circle, and its exact defect}
\author{}
\date{23 October 2026}
\begin{document}
\maketitle
\vspace{-1.5em}

This note replaces the residue-selected Bob parameter in the leading
unit-$L$ form by six fixed functions.  An explicit source-only
perturbation bound then gives a finite fixed testing family for the
complete actual C1 parent on a complementary annulus.  No certificate
on the full testing domain is asserted: the last proposition identifies
the negative term obstructing a three-event cap at the omitted corner.
The source conventions are those of
\texttt{pda4a21937f64}, \texttt{pb0d9d28d0849}, and
\texttt{p23a094ff816d} in commit
\texttt{fb58c232e374dadf3cbb5f23e7e7413d429de884}.

\section*{The six-function completion}

Let
\[
 L>0,\qquad 0<r<1,\qquad Lr<1,
 \qquad
 C_0=\frac1{2L(1-r)},\quad
 A_0=\frac{1+L(1-r)}{2Lr},\quad
 a_0=\frac L{2(1-Lr)}.
\]
The inequalities on $r$ and $Lr$ are strict.  All limits towards their
openings are allowed; no value is assigned to a divergent coefficient
at a literal zero denominator.  Write
\[
 B=\frac{1+L(1-r)}{2Lr(1-r)},\qquad
 \gamma=\frac{1+L(1-r)}{2L}
       =\frac12+\frac1{4a_0}.
\]
Let $c^2+s^2=1$, with $c\geq0$.  The endpoints $c=0$, $s=\pm1$
are included.  The four state variables $P,X,Y,T$ below are independent.
For finite physical angles the substitution is
\[
 c=(1+v^2)^{-1/2},\quad s=vc,\quad P=uW,\quad T=cZ.
\]
Identities on the independent raw variables therefore retain the full
state polynomial; they impose neither a state graph nor a relation
between the two state normals.

\begin{lemma}[Residue-free Bob completion]
For every source in the displayed domain,
\begin{align}
 M(s)&=(C_0+c^2/2)(P-Y)^2+(A_0+s^2/2)Y^2
             +a_0(sY-T)^2-sPT\notag\\
 &=a_0\left(T-sY-\frac{1-Lr}{L}sP\right)^2
        +B(Y-rP)^2+\gamma(cP)^2.                 \label{eq:three}
\end{align}
Thus the fixed six-function family
\[
 (T,sY,sP,Y,P,cP)
\]
has the positive source Gram matrix
\[
 \begin{pmatrix}
 a_0&-a_0&-1/2&0&0&0\\
 -a_0&a_0&1/2&0&0&0\\
 -1/2&1/2&1/(4a_0)&0&0&0\\
 0&0&0&B&-Br&0\\
 0&0&0&-Br&Br^2&0\\
 0&0&0&0&0&\gamma
 \end{pmatrix}.                                    \label{eq:sixgram}
\]
It has rank three, including at either Bob endpoint as an abstract
feature Gram.  Evaluation can lower the rank because the functions
can become dependent there.
\end{lemma}

\begin{proof}
Put $k=C_0+1/2$.  The source identities
\[
 C_0+A_0+1/2=B,\qquad k/B=r,\qquad kA_0/B=\gamma
\]
reduce the claim to two successive completions of squares, first in
$T$ and then in $Y$.  They also give the three rank-one summands of
\eqref{eq:sixgram}.  No source or testing residue is chosen.
\end{proof}

\section*{A complete four-state annulus}

The normalized complete signed core, in the same raw variables, is
\begin{align}
 K={}&X^2+c^2(P-Y)^2+s^2Y^2-2sPT
       +2ucX(sP-sY+T)\notag\\
    &+u^2s^2X^2+u^2(sY-T)^2.                    \label{eq:core}
\end{align}
Under $P=uW$, $T=cZ$, this is identically
\[
 \{c(P-Y)+usX-sZ\}^2
 +(1+u^2)(sY-cZ)^2+(X+uZ)^2-(1+u^2)Z^2.
\]
The whole leading selected quarter-parent is
\[
 S_0=C_0(P-Y)^2+A_0(1+u^2)Y^2
           +a_0(1+u^2)(sY-T)^2+K/2.              \label{eq:parent}
\]
In particular, the independent $X$ normal and all $u$-dependent terms
remain present.

\begin{proposition}[Explicit uniform leading-source reserve]
Fix $0<\ell\leq U$, $0<\tau\leq1$.  If
\[
 \ell\leq L\leq U,\quad |u|\leq1,\quad
 c^2+s^2=1,\quad c\geq0,\quad u^2+c^2\geq\tau^2,
\]
the matrix of $S_0$ in $(P,X,Y,T)$ satisfies
\[
 G\succeq \kappa I_4,\qquad
 \kappa=\frac{\min(\ell,1)\tau^2}
                   {4(7U^2+5U+3)}.               \label{eq:kappa}
\]
This estimate is uniform through the three source openings
$r\to0$, $r\to1$, and $Lr\to1$ whenever they are compatible with
the displayed source domain.
\end{proposition}

\begin{proof}
The restriction $T=0$ is the strict three-variable sum of squares
\begin{align*}
 &C_0(P-Y)^2+A_0(1+u^2)Y^2+a_0(1+u^2)s^2Y^2\\
 &\hspace{1em}+\tfrac12\{c(P-Y)+usX\}^2
       +\tfrac12X^2+\tfrac12s^2(1+u^2)Y^2.
\end{align*}
The exact determinant is
\[
 \det G=\frac{(1+u^2)^2\{Lc^2+u^2s^2\}
                    \{1+L(1-r)\}^2}
                   {16L^2r(1-r)(1-Lr)}.
\]
It is positive unless $u=c=0$.  The strict upper state block and a
one-field Schur complement therefore prove $G\succ0$.  This uses a
complete state matrix, not a determinant without a normal-block check.

For the quantitative estimate put $t=u^2$, $q=c^2$, $e=1-r$ and
$D=Lq+t(1-q)$.  Direct cofactors, with all opening denominators
cancelled, give
\[
 \operatorname{tr}G^{-1}
 =\frac{2(L^2N_2+LN_1+N_0)}{(1+t)D(1+Le)},
\]
where
\begin{align*}
 N_2&=1+r^2+eq\{1+r(1-q)\}+e(2-qr)t+e^2(1-q)t^2,\\
 N_1&=2eq+\{2-rq(3-q)\}t+e(1-q)t^2,\\
 N_0&=1+(2-q)t.
\end{align*}
For $0\leq q,t\leq1$, $0<r<1$, these satisfy
$0\leq N_2\leq7$, $0\leq N_1\leq5$, $0\leq N_0\leq3$.
Also
\[
 D\geq\min(\ell,1)\{q+t(1-q)\}
    \geq\tfrac12\min(\ell,1)(q+t).
\]
Thus $\operatorname{tr}G^{-1}\leq1/\kappa$.
The largest eigenvalue of $G^{-1}$ is at most its trace, proving
\eqref{eq:kappa}.
\end{proof}

\section*{Explicit fixed testing functions on that annulus}

The same reserve admits actual source weights on an explicit domain.
For positive $C,A,a$ put
\[
 d_{\rm src}=|C/C_0-1|+|A/A_0-1|+|a/a_0-1|.
\]
If the leading raw Gram is at least $kI$, the actual selected raw Gram
is at least $\{k-(k+3)d_{\rm src}\}I$.  Indeed put $f=1-d_{\rm src}$,
retain the nonnegative event excesses $C-fC_0$, $A-fA_0$, $a-fa_0$,
and use the exact identity
\[
 S(C,A,a)=fS_0+(1-f)K/2+\text{the three event excesses}.
\]
The raw matrix of $K/2$ has absolute row sums at most $5/2$, hence
operator norm at most $3$.  This proves the bound with both state
normals retained.

Here is a finite construction, rather than an appeal to pointwise
bounded-Gram positivity.  Choose rational $\ell,U,\tau$ in the displayed
ranges, so that every fixed coefficient below is rational or algebraic.
Set
\[
 k_* =\frac{\min(\ell,1)\tau^2}{8(7U^2+5U+3)},\quad
 N=4+\frac{216}{k_*},\quad h_* =\min(1,k_*/24),\quad B_*=N+1.
\]
The actual source admission is
\[
 C,A,a>0,\qquad
 d_{\rm src}\leq\frac{k_*}{2(k_*+3)}.             \label{eq:actual-admission}
\]
It implies an actual raw reserve $k_*/2$ on the slightly larger annulus
$u^2+c^2\geq\tau^2/2$.  All infinitesimal unit-$L$ sources with
$\ell\leq L\leq U$ satisfy this admission, by the unchanged opening-safe
ratios LA1/BE9.  The finite statement here uses the displayed inequality
itself and does not infer an explicit $x,\sigma$ bound from an
asymptotic assertion.

Clip the three source weights at $N$:
$C'=\min(C,N)$, $A'=\min(A,N)$,
$a'=\min(a,N)$.  Retain the three nonnegative differences.
Use the event coordinates
\[
 e=(P-Y,Y,sY-T,X)^T.
\]
In these coordinates the clipped Gram has the exact form
\[
 G'=G_{\rm const}+\sum_{j=1}^7\eta_j A_j,
 \qquad G_{\rm const}=\diag(C',A',a',1/2),
\]
with testing monomials
\[
 \eta=(u^2,c^2,s^2,s,ucs,uc,u^2s^2).
\]
The matrices, with indices $1,\ldots,4$ and $E_{ij}$ denoting matrix
units, are
\begin{align*}
 A_1&=\diag(0,A',a'+1/2,0),&A_2&=E_{11}/2,\\
 A_3&=-(E_{12}+E_{21}+E_{22})/2,
 &A_4&=(E_{13}+E_{31}+E_{23}+E_{32})/2,\\
 A_5&=(E_{14}+E_{41}+E_{24}+E_{42})/2,
 &A_6&=-(E_{34}+E_{43})/2,\\
 A_7&=E_{44}/2.
\end{align*}
All satisfy $\|A_j\|\leq B_*$.

For clarity, the uniform clipping claim used here has a short proof.
The core matrix in the event coordinates has row sums at most $5/2$,
so its operator norm is at most $3$.  The inverse change of coordinates
has squared Frobenius norm at most $6$.  At any point with
$u^2+c^2\geq\tau^2/2$, the preceding proposition therefore gives the
unclipped actual event-coordinate Gram at least $(k_*/12)I$.
Split coordinates according as their source event weight exceeds $N$.
The unchanged complementary principal block is at least $k_*/12$.
The clipped large block is at least $N-3$, and the off-diagonal
block has norm at most $3$.  The inequality
\[
 6ab\leq\frac{216}{k_*}a^2+\frac{k_*}{24}b^2
\]
proves $G'\succeq h_*I$.  This argument uses only source clipping;
the indices being clipped do not depend on a test.

Choose a fixed integer
\[
 n\geq\max\left(\frac8{\tau^2},\frac{56B_*}{h_*}\right),
 \qquad \delta=4/n.
\]
Take the finite algebraic grid
\[
 u_j=j/n,\qquad
 c_k=\cos\frac{k\pi}{2n},\qquad
 s_k=\sin\frac{k\pi}{2n},\qquad -n\leq j,k\leq n,
\]
retaining centers $a=(u_j,c_k,s_k)$ with
$u_j^2+c_k^2\geq\tau^2/2$.  A chart consists of circle points within
$1/n$ in each coordinate of its center.  These charts cover the
target annulus: nearest-grid errors are at most $1/(2n)$ in $u$
and $\pi/(4n)<1/n$ in either circle coordinate, and the change in
$u^2+c^2$ is at most $4/n$.

On a fixed chart put $b_j=\eta_j-\eta_j(a)$.
Then $|b_j|\leq\delta$.  The exact matrix identity
\begin{align}
 G'={}&\left(G'(a)-7\delta B_* I\right)\notag\\
 &+\frac12\sum_{j=1}^7
  \left[(B_*I+A_j)(\delta+b_j)
          +(B_*I-A_j)(\delta-b_j)\right]          \label{eq:fixed}
\end{align}
has only positive source matrices: its first block is at least
$h_*/2$, and $B_*I\pm A_j\succeq0$.  Consequently the following
\emph{fixed}, source-independent testing functions suffice on that chart:
\[
 e_i,\quad e_i\sqrt{\delta+b_j},\quad
 e_i\sqrt{\delta-b_j}\quad(1\leq i\leq4,\ 1\leq j\leq7),
\]
together with
\[
 \sqrt{1+u^2}\,Y,\qquad
 \sqrt{1+u^2}\,(sY-T).
\]
There are at most $62$ functions per chart, and at most
$(2n+1)^2$ charts.  The excess $C-C'$ uses the already retained
function $e_1$; the other two excesses use the last two functions.
Their source Gram coefficients are $C-C'$, $A-A'$ and $a-a'$.
The blocks in \eqref{eq:fixed}, the displayed $A_j$, and these excess
coefficients are the complete evaluated matrix prescription.  All
centers are algebraic, all source entries are rational on the eight
possible source clipping chambers, and every radical has a fixed
constant $\delta$.  Factoring the positive matrices takes place only
in the source field.  These identities also hold over separately
labelled real closed fields, chart by chart.

For the actual C1 source put $a=\alpha$.  After the physical substitution
$P=uW$, $T=cZ$, restore the three fixed testing functions
\[
 r_m=sW-cX+usY-ucZ,\quad
 r_D=c(uW-Y)+usX-sZ,\quad r_\beta=X+uZ.
\]
Their source coefficients are respectively $m,D,\beta\geq0$.
The complete normalized parent, with $\lambda_{\rm ev}=n=0$, is
identically $S(C,A,\alpha)+mr_m^2+Dr_D^2+\beta r_\beta^2$.
Thus at most $65$ fixed functions per chart admit the \emph{entire}
actual C1 parent on this annulus.  All $m,D,\beta\geq0$ are allowed,
including literal zero coefficients.  All state variables and the
literal Bob endpoints remain included.  The restriction $|u|\leq1$
and the omitted double corner remain part of the theorem's domain.
Real and imaginary splitting handles complex states without new
testing functions.

\section*{The exact obstruction for actual stationary weights}

\begin{proposition}[Below-wall defect and a full selected-parent witness]
Let $C,A,a>0$ be arbitrary and set
\[
 k=C+1/2,\quad B=C+A+1/2,\quad
 g=\frac{kA}{B},\quad \gamma=\frac12+\frac1{4a},
\]
\[
 \Delta=8CAa-4Ca-1-2(C+A+a),\qquad
 d_*=\frac{\Delta}{8aB}=g-\gamma.
\]
The moving Bob form obeys the exact identity
\begin{align}
 &(C+c^2/2)(P-Y)^2+(A+s^2/2)Y^2+a(sY-T)^2-sPT\notag\\
 &\quad=a\left(T-sY-\frac{sP}{2a}\right)^2
       +B\left(Y-\frac{k}{B}P\right)^2
       +g c^2P^2+d_*s^2P^2.                    \label{eq:defect}
\end{align}
The leading weights have $\Delta=0$.  If $\Delta<0$, the selected
complete parent with weights $(C,A,a)$ is negative at an explicit
finite angle pair and a finite nonzero state.
\end{proposition}

\begin{proof}
Two square completions prove \eqref{eq:defect}; the coefficient identity
$8aB(g-\gamma)=\Delta$ is immediate by multiplication.
For the last assertion retain the complete core \eqref{eq:core} and put
\[
 X=0,\quad P=1,\quad Y=k/B,\quad
 T=sY+s/(2a),\quad
 E=A(k/B)^2+F>0,\qquad F=\frac{a+1/2}{4a^2}>0.
\]
Its complete selected value is exactly
\[
 d_*+\gamma c^2+Eu^2-Fu^2c^2.
\]
Choose
\[
 t=\frac{-d_*}{2(\gamma+E+1)},\quad
 u=c=\sqrt t,\quad s=\sqrt{1-t}.
\]
Since $g>0$, one has $0<t<1/2$.  The selected value becomes
\[
 d_*\frac{\gamma+E+2}{2(\gamma+E+1)}-Ft^2<0.
\]
The actual finite angle and state realization is
$v=s/c$, $W=1/u$, $X=0$, $Y=k/B$, $Z=T/c$.
Thus the witness uses the entire selected four-state polynomial,
not a determinant or a state-normal truncation.
\end{proof}

For an actual finite noncommuting stationary contact of the full
nonnegative-event parent, the existing strict three-event result
(CAP14) gives $\Delta(C,A,\alpha)<0$.  Substituting $a=\alpha$
in \eqref{eq:defect} therefore exposes a strictly negative term which
the leading source loses.  The explicit witness disproves replacing
the full source by that three-event cap.  It is not a counterexample
to the full parent: the omitted original events supply indispensable
contact-scale contributions.  The source identities BE3--BE9, the
$e$-bound and SP6 have neither been altered nor replaced.  Literal
$D=0$ or $\beta=0$ alone do not remove the defect argument.

The annulus construction absorbs this defect only where the displayed
strict reserve is available.  The missing global step remains a fixed
separated completion with the other original events at the double
contact corner, projective Alice charts, and the remaining source
scales and regions.

\end{document}
